% area: Stochastic Calculus

\begin{definition}[Sigma-Algebra]
\label{def:sigma-algebra}
Let $\Omega$ be a nonempty set. A \emph{sigma-algebra} (or $\sigma$-algebra) on $\Omega$ is a collection $\mathcal{F} \subseteq 2^\Omega$ of subsets satisfying:
\begin{enumerate}
  \item $\Omega \in \mathcal{F}$.
  \item If $A \in \mathcal{F}$, then $A^c \in \mathcal{F}$ (closed under complements).
  \item If $A_1, A_2, \ldots \in \mathcal{F}$, then $\bigcup_{n=1}^\infty A_n \in \mathcal{F}$ (closed under countable unions).
\end{enumerate}
The pair $(\Omega, \mathcal{F})$ is called a \emph{measurable space}, and elements of $\mathcal{F}$ are called \emph{measurable sets} or \emph{events}.
\end{definition}

\begin{definition}[Probability Space]
\label{def:prob-space}
A \emph{probability space} is a triple $(\Omega, \mathcal{F}, \mathbb{P})$ where:
\begin{enumerate}
  \item $(\Omega, \mathcal{F})$ is a measurable space (Definition~\ref{def:sigma-algebra}), with $\Omega$ the \emph{sample space}.
  \item $\mathbb{P} : \mathcal{F} \to [0,1]$ is a \emph{probability measure}: $\mathbb{P}(\Omega) = 1$, and for pairwise disjoint $A_1, A_2, \ldots \in \mathcal{F}$,
  \[
    \mathbb{P}\!\left(\bigcup_{n=1}^\infty A_n\right) = \sum_{n=1}^\infty \mathbb{P}(A_n).
  \]
\end{enumerate}
\end{definition}

\begin{definition}[Random Variable]
\label{def:random-variable}
Let $(\Omega, \mathcal{F}, \mathbb{P})$ be a probability space (Definition~\ref{def:prob-space}).
A \emph{random variable} is a measurable function $X : \Omega \to \mathbb{R}$, meaning $X^{-1}(B) \in \mathcal{F}$ for every Borel set $B \subseteq \mathbb{R}$.
The \emph{law} (or \emph{distribution}) of $X$ is the pushforward measure $\mathbb{P}_X(B) := \mathbb{P}(X^{-1}(B))$ on $\mathbb{R}$.
\end{definition}

\begin{definition}[Expectation]
\label{def:expectation}
Let $X$ be a random variable on $(\Omega, \mathcal{F}, \mathbb{P})$ (Definition~\ref{def:random-variable}).
The \emph{expectation} of $X$ is
\[
  \mathbb{E}[X] := \int_\Omega X(\omega)\, d\mathbb{P}(\omega),
\]
provided the integral exists. For $p \geq 1$, we say $X \in L^p(\Omega, \mathcal{F}, \mathbb{P})$ if $\mathbb{E}[|X|^p] < \infty$.
\end{definition}

\begin{definition}[Independence]
\label{def:independence}
Let $(\Omega, \mathcal{F}, \mathbb{P})$ be a probability space (Definition~\ref{def:prob-space}).
\begin{enumerate}
  \item Events $A, B \in \mathcal{F}$ are \emph{independent} if $\mathbb{P}(A \cap B) = \mathbb{P}(A)\mathbb{P}(B)$.
  \item Random variables $X, Y$ (Definition~\ref{def:random-variable}) are \emph{independent} if $\sigma(X)$ and $\sigma(Y)$ are independent sub-$\sigma$-algebras, i.e.\ $\mathbb{P}(X \in A, Y \in B) = \mathbb{P}(X \in A)\mathbb{P}(Y \in B)$ for all Borel $A, B$.
  \item A collection $\{X_i\}_{i \in I}$ is \emph{independent} if for every finite $J \subseteq I$, $\mathbb{P}(\bigcap_{j \in J} \{X_j \in B_j\}) = \prod_{j \in J} \mathbb{P}(X_j \in B_j)$.
\end{enumerate}
\end{definition}

\begin{definition}[Conditional Expectation]
\label{def:conditional-expectation}
Let $X \in L^1(\Omega, \mathcal{F}, \mathbb{P})$ (Definition~\ref{def:expectation}) and let $\mathcal{G} \subseteq \mathcal{F}$ be a sub-$\sigma$-algebra (Definition~\ref{def:sigma-algebra}).
The \emph{conditional expectation} $\mathbb{E}[X \mid \mathcal{G}]$ is the $\mathbb{P}$-a.s.\ unique $\mathcal{G}$-measurable random variable satisfying
\[
  \int_G \mathbb{E}[X \mid \mathcal{G}]\, d\mathbb{P} = \int_G X\, d\mathbb{P} \quad \text{for all } G \in \mathcal{G}.
\]
Its existence and uniqueness follow from the Radon--Nikod\'{y}m theorem.
\end{definition}

\begin{definition}[Borel Sigma-Algebra]
\label{def:borel}
The \emph{Borel $\sigma$-algebra} $\mathcal{B}(\mathbb{R})$ on $\mathbb{R}$ is the smallest $\sigma$-algebra (Definition~\ref{def:sigma-algebra}) containing all open subsets of $\mathbb{R}$.
More generally, for a topological space $X$, the Borel $\sigma$-algebra $\mathcal{B}(X)$ is generated by the open sets of $X$.
\end{definition}

\begin{definition}[Measure Space]
\label{def:measure-space}
A \emph{measure space} is a triple $(\Omega, \mathcal{F}, \mu)$ where $(\Omega, \mathcal{F})$ is a measurable space (Definition~\ref{def:sigma-algebra}) and $\mu : \mathcal{F} \to [0, \infty]$ satisfies $\mu(\emptyset) = 0$ and countable additivity.
It is \emph{$\sigma$-finite} if $\Omega = \bigcup_{n=1}^\infty A_n$ with $\mu(A_n) < \infty$ for each $n$.
A probability space (Definition~\ref{def:prob-space}) is a measure space with $\mu(\Omega) = 1$.
\end{definition}

\begin{theorem}[Monotone Convergence Theorem]
\label{thm:mct}
Let $(\Omega, \mathcal{F}, \mu)$ be a measure space (Definition~\ref{def:measure-space}) and let $0 \leq f_1 \leq f_2 \leq \cdots$ be a non-decreasing sequence of non-negative measurable functions with $f_n \to f$ pointwise.
Then
\[
  \lim_{n \to \infty} \int f_n\, d\mu = \int f\, d\mu.
\]
\end{theorem}

\begin{theorem}[Dominated Convergence Theorem]
\label{thm:dct}
Let $(\Omega, \mathcal{F}, \mu)$ be a measure space (Definition~\ref{def:measure-space}).
Suppose $f_n \to f$ pointwise $\mu$-a.e.\ and $|f_n| \leq g$ $\mu$-a.e.\ for all $n$, where $g \in L^1(\mu)$.
Then $f \in L^1(\mu)$ and
\[
  \lim_{n \to \infty} \int f_n\, d\mu = \int f\, d\mu.
\]
\end{theorem}

\begin{theorem}[Radon--Nikod\'{y}m Theorem]
\label{thm:radon-nikodym}
Let $(\Omega, \mathcal{F})$ be a measurable space (Definition~\ref{def:sigma-algebra}) and let $\mu, \nu$ be $\sigma$-finite measures (Definition~\ref{def:measure-space}) with $\nu \ll \mu$ (i.e.\ $\mu(A) = 0 \Rightarrow \nu(A) = 0$).
Then there exists a non-negative measurable function $f$, called the \emph{Radon--Nikod\'{y}m derivative} $\frac{d\nu}{d\mu}$, such that
\[
  \nu(A) = \int_A f\, d\mu \quad \text{for all } A \in \mathcal{F}.
\]
The function $f$ is unique $\mu$-a.e.
\end{theorem}

\begin{definition}[Stochastic Process]
\label{def:stochastic-process}
Let $(\Omega, \mathcal{F}, \mathbb{P})$ be a probability space (Definition~\ref{def:prob-space}) and let $T \subseteq [0, \infty)$.
A \emph{stochastic process} indexed by $T$ is a collection $\{X_t\}_{t \in T}$ of random variables (Definition~\ref{def:random-variable}) on $\Omega$.
For each fixed $\omega \in \Omega$, the map $t \mapsto X_t(\omega)$ is called a \emph{sample path} (or \emph{trajectory}) of the process.
\end{definition}

\begin{definition}[Filtration and Adapted Process]
\label{def:filtration}
A \emph{filtration} on $(\Omega, \mathcal{F}, \mathbb{P})$ (Definition~\ref{def:prob-space}) is an increasing family $(\mathcal{F}_t)_{t \geq 0}$ of sub-$\sigma$-algebras of $\mathcal{F}$:
\[
  s \leq t \implies \mathcal{F}_s \subseteq \mathcal{F}_t \subseteq \mathcal{F}.
\]
A filtration is \emph{right-continuous} if $\mathcal{F}_t = \bigcap_{s > t} \mathcal{F}_s$ for all $t$.
A process $\{X_t\}$ (Definition~\ref{def:stochastic-process}) is \emph{adapted} to $(\mathcal{F}_t)$ if $X_t$ is $\mathcal{F}_t$-measurable for every $t \geq 0$.
The \emph{natural filtration} of $\{X_t\}$ is $\mathcal{F}_t^X := \sigma(X_s : s \leq t)$.
\end{definition}

\begin{definition}[Martingale]
\label{def:martingale}
Let $\{X_t\}_{t \geq 0}$ be an adapted process (Definition~\ref{def:filtration}) with $X_t \in L^1$ (Definition~\ref{def:expectation}) for all $t$.
The process is a \emph{martingale} with respect to $(\mathcal{F}_t)$ if
\[
  \mathbb{E}[X_t \mid \mathcal{F}_s] = X_s \quad \mathbb{P}\text{-a.s.\ for all } 0 \leq s \leq t.
\]
It is a \emph{supermartingale} (resp.\ \emph{submartingale}) if the equality is replaced by $\leq$ (resp.\ $\geq$).
\end{definition}

\begin{definition}[Stopping Time]
\label{def:stopping-time}
Let $(\mathcal{F}_t)_{t \geq 0}$ be a filtration (Definition~\ref{def:filtration}).
A random variable $\tau : \Omega \to [0, \infty]$ is a \emph{stopping time} with respect to $(\mathcal{F}_t)$ if
\[
  \{\tau \leq t\} \in \mathcal{F}_t \quad \text{for all } t \geq 0.
\]
The \emph{stopped $\sigma$-algebra} is $\mathcal{F}_\tau := \{ A \in \mathcal{F} : A \cap \{\tau \leq t\} \in \mathcal{F}_t \text{ for all } t \geq 0 \}$.
\end{definition}

\begin{definition}[Standard Brownian Motion]
\label{def:brownian-motion}
A \emph{standard Brownian motion} (or \emph{Wiener process}) is a stochastic process $\{B_t\}_{t \geq 0}$ (Definition~\ref{def:stochastic-process}) on a probability space $(\Omega, \mathcal{F}, \mathbb{P})$ (Definition~\ref{def:prob-space}) satisfying:
\begin{enumerate}
  \item \textbf{Initial value:} $B_0 = 0$ $\mathbb{P}$-a.s.
  \item \textbf{Independent increments:} For $0 \leq t_0 < t_1 < \cdots < t_n$, the increments $B_{t_1} - B_{t_0}, \ldots, B_{t_n} - B_{t_{n-1}}$ are independent (Definition~\ref{def:independence}).
  \item \textbf{Stationary Gaussian increments:} For $s < t$, $B_t - B_s \sim \mathcal{N}(0, t-s)$.
  \item \textbf{Continuous paths:} $t \mapsto B_t(\omega)$ is continuous $\mathbb{P}$-a.s.
\end{enumerate}
\end{definition}
